\chapter{PartDesign: turning sketches into solids}

PartDesign is MobileCAD's most-used workbench. Its core idea is \textbf{history}: a body (Body) is built from a chain of
features, and each feature does one thing to the result of the previous step. Change any step and the steps after it
recompute automatically.

\section{Three key ideas: body, feature, Tip}

\begin{description}
  \item[Body] The container for features, and the unit of work in PartDesign. Every sketch and feature belongs to a body.
        When you create a feature and there is no body yet, the app creates one and puts the new feature in it.
  \item[Feature] One step in the history: pad, pocket, fillet, array, and so on.
  \item[Tip] The marker for ``how far the body has been computed''. A blue rollback bar runs through the tree; long-press
        a feature and choose \ui{Roll tip to here} to return to the state at that step (replaying the modeling process).
\end{description}

\textbf{Active body} (Active Body): with several bodies, new features go into the ``active'' one (its label is bold in the
tree). Double-tap a body, or right-click \ui{Toggle active body}, to switch.

\section{How a feature gets prepared}

Most PartDesign commands read the selection and act on it, which saves a lot of dialog work:

\begin{itemize}
  \item Select a \textbf{sketch} and tap Pad or Pocket → the profile is filled in;
  \item Select a \textbf{plane} and tap Pad or Pocket → a new sketch is created on that plane first;
  \item With nothing selected, Pad or Pocket uses the \textbf{last sketch} in the document; if there is none, it says
        ``select a sketch or a plane first'';
  \item Fillet/chamfer/draft/thickness: if you select a whole feature, they apply to all edges or faces automatically
        (\ui{UseAllEdges}); if nothing is selected, they wait at the Tip for you to tap edges;
  \item Defaults: Pad 10 mm, Pocket 5 mm, Revolution 360° (axis: the sketch V axis), Mirror on the V axis, LinearPattern
        100 mm × 2 along X, PolarPattern 360° × 2 about the Z axis.
\end{itemize}

\section{Additive features}

\small
\begin{longtable}{@{}p{0.17\linewidth}p{0.10\linewidth}p{0.65\linewidth}@{}}
\toprule \textbf{Feature} & \textbf{Command} & \textbf{Key parameters} \\ \midrule \endhead
Pad & Pad & Profile, type (length/two lengths/to last/to face/to shape/dimension), length, length 2, offset, midplane, reverse, draft angle, refine \\
Revolution & Revolution & Profile, reference axis, type (angle/two angles/to last/to face), angle, angle 2, midplane, reverse, refine \\
Loft & Additive Loft & Profile, sections (several), ruled, closed, refine \\
Pipe / Sweep & Additive Pipe & Profile, path, mode, transition, transformation, refine \\
Helix & Additive Helix & Profile, reference axis, mode, pitch, height, turns, angle, growth, left-handed, reverse, refine \\
Additive primitives & Additive Box/Cylinder/Sphere/Cone/Torus/Ellipsoid/Prism/Wedge & Length, width, height/radius/angle and so on; created on the face when one is selected \\
\bottomrule
\end{longtable}

\section{Subtractive features}

\begin{itemize}
  \item \ui{Pocket}: the mirror image of Pad --- it scoops the profile inward. The depth types likewise include through all;
  \item \ui{Hole}: the standard hole feature. Parameters: profile, \ui{Thread}, thread type/size/class/fit, diameter,
        direction, hole cut type (counterbore/countersink/counterdrill) with its diameter/depth/angle, depth type/depth,
        drill point angle, taper, reverse;
  \item \ui{Groove}: cuts an annular groove around an axis;
  \item \ui{Subtractive loft/pipe/helix} take the same parameters as the additive versions, only in the opposite direction;
  \item \ui{Subtractive primitives}: the subtractive versions of the 8 primitives.
\end{itemize}

\begin{tip}[Hole finds its own position]
Call Hole on a face with no circular hole and the app builds an R3 sketch circle at the center of that face as the
profile --- then change the hole diameter. This saves the ``draw a circle first'' step.
\end{tip}

\section{Dress-up features}

\begin{description}
  \item[Fillet] Parameters: base, edges (an edge table, with a radius per edge), use all edges, refine. Drag the arrow to change the radius directly on the model.
  \item[Chamfer] Parameters: base, chamfer type (equal distance/two distances/distance + angle), size, size 2, angle, flip direction, use all edges.
  \item[Draft] Parameters: base, neutral plane, pull direction, angle, reverse.
  \item[Thickness] Parameters: base, thickness value, mode (skin/pipe/rectangular), join type (arc/tangent/intersection), reverse, intersection.
  \item[Defeaturing] Removes existing faces or fillets and restores the geometry underneath.
\end{description}

\screenshot{ui-feature-dialog-zh.jpg}{A feature dialog (cone): the main properties come first (angle, height, radius 1, radius 2), the rest fold into ``More options (6)''. The ✓/✗ in the top-right corner commits or rolls back.}

\section{Arrays and mirrors}

\begin{center}
\small
\begin{tabularx}{\linewidth}{@{}lX@{}}
\toprule
\textbf{Feature} & \textbf{Parameters} \\
\midrule
LinearPattern & Original feature, direction (edge/axis), mode (extent/spacing), length, offset, count, reverse, refine \\
PolarPattern & Original feature, axis, mode (angle/spacing), angle, offset, count, reverse, refine \\
Mirrored & Original feature, mirror plane, refine \\
MultiTransform & Combines several arrays/mirrors (uses the generic object-creation flow)\\
\bottomrule
\end{tabularx}
\end{center}

\begin{tip}[Extent vs spacing in an array]
\textbf{Extent} is the total distance from the first instance to the last; \textbf{spacing} is the distance between two
neighbouring instances. For six holes around a circle: choose \ui{PolarPattern}, the Z axis, an angle of 360° and a
count of 6.
\end{tip}

\section{Datum geometry and attachment}

\ui{Datum} geometry gives you a ``virtual reference'' when no existing face will do:

\begin{itemize}
  \item \ui{Datum plane}, \ui{Datum line}, \ui{Datum point}, \ui{Coordinate system};
  \item Parameters at creation: attachment support (\ui{AttachmentSupport}), map mode (\ui{MapMode}), \ui{Reverse}, offset
        (\ui{AttachmentOffset}), path parameter (\ui{MapPathParameter});
  \item \ui{Edit attachment} (long-press in the tree) changes only the attachment and does not open the whole feature.
\end{itemize}

\begin{warn}[The attachment mode decides what follows what]
When a sketch hangs on a face of a body and that face moves because an earlier feature changed, the sketch and the
features derived from it update with it --- exactly what parametric modeling is for. But if the face disappears (cut away
by a chamfer, for example), the features that depend on it fail.
\end{warn}

\section{Boolean operations}

\begin{itemize}
  \item \ui{PartDesign → Boolean}: uses the selected objects as tools and fuses, cuts or intersects them with the active body;
  \item The \ui{Fuse/Cut/Common} commands of the \ui{Part} workbench are Booleans between independent solids, not part of
        the PartDesign history;
  \item Before a Boolean, check that both shapes are valid solids (\ui{Part → Check geometry}).
\end{itemize}

\section{Reading a feature dialog}

A feature dialog has three parts:

\begin{enumerate}
  \item \textbf{Title bar}: name + ▾ collapse + ✓ OK + ✗ Cancel;
  \item \textbf{Error bar} (red): shows the reason when the feature fails, for example ``the geometry does not exist on
        the base feature'';
  \item \textbf{Property area}: the main properties in FreeCAD task-panel order, the rest under \ui{More options (N)}.
\end{enumerate}

How each kind of property row works is described in section 3.7. Two practical details:

\begin{itemize}
  \item \textbf{Live preview}: a change acts on the model at once, and each change is its own transaction;
  \item \textbf{Dragging the arrow}: quantity properties such as pad length or fillet radius have an arrow in the view;
        drag it and read the value in millimetres on the label.
\end{itemize}

\section{Failed features and diagnosis}

When a feature fails:

\begin{itemize}
  \item Its label turns \textbf{red} in the tree and a red error icon appears;
  \item The view shows its last valid predecessor shape as a \textbf{red translucent ``ghost''};
  \item Open the dialog: the red bar at the top gives the error string.
\end{itemize}

\begin{practice}[A four-step diagnosis]
1) Open the failed feature and read the red bar; 2) check whether the profile, face or edge it references still exists; 3) long-press it and choose \ui{Suppress} to see whether the features after it depend on it too;
4) go back to the sketch and fix the constraints, or long-press and choose \ui{Recompute object}.
\end{practice}

\section{A complete workflow: from a block to a finished part}

Continue with the bracket from chapter 4:

\begin{enumerate}
  \item \textbf{Add a rib}: create a sketch on the top face of the bracket, draw a line → leave the sketch → \ui{Pad}
        (tick ``symmetric to plane'') to get the rib plate;
  \item \textbf{Drill a countersunk hole}: select the top face → \ui{Hole} → turn on \ui{Thread} or select
        \ui{Countersink} → enter diameter/depth;
  \item \textbf{Array}: select the hole → \ui{LinearPattern}, pick an edge for the direction, count 4, length 40;
  \item \textbf{Chamfer}: select the outer contour edges → \ui{Chamfer}, size 1 mm;
  \item \textbf{Change the base}: double-tap the first Pad and change the length from 12 to 16 --- the rib, holes, array
        and chamfer all follow automatically;
  \item Export STEP for a colleague, or go to the ``Drawing'' tab and make a drawing.
\end{enumerate}

\screenshot[0.72]{tablet-partdesign.jpg}{A real PartDesign model (a screwdriver bit): the tree reads Pad → Groove → Mirrored → Revolution → DatumPlane → Pocket → PolarPattern... Every step can still be changed, which is exactly the power of parametric modeling.}

\section{Common pitfalls}

\begin{itemize}
  \item \textbf{A sketch must be inside the active body}; a cross-body reference is replaced by a SubShapeBinder, which
        can raise a circular-dependency error;
  \item \textbf{There is no M key}: desktop shortcuts such as switching polyline mode or cycling dimension types do not
        exist on mobile;
  \item \textbf{Polygon side counts are fixed}: the commands offer 3--8 sides only;
  \item \textbf{Multi-point tools need a double-tap or Finish}: a polyline or B-spline is not committed until you do;
  \item \textbf{Some arrays use the generic flow}: circular/path/point arrays and MultiTransform, and the main parameters
        of ellipsoid/prism/wedge, are under ``More options'';
  \item \textbf{Validate Sketch shows no repair count}: it does repair, but the interface reports no number.
\end{itemize}
